% part: intuitionistic-logic
% chapter: propositions-as-types
% section: proof-terms

\documentclass[../../../include/open-logic-section]{subfiles}

\begin{document}

\olfileid{int}{pty}{ter}

\olsection{Proof Terms}

We label assumptions in \Log{N1i} with variables (e.g., $!A^x$) as
discharge labels. In \Log{N2i}, every sequent $\Gamma \Sequent !A$
has a list of variable--!!{formula} pairs in the context~$\Gamma$ on
the left. Thus, in \Log{N2i}, every sequent in !!a{proof} carries not
only the information what the !!{proof} so far has !!{prove}d (i.e.,
$!A$), but also which assumptions are open at each point (i.e.,
$\Gamma$). What if we also included in each sequent the information
\emph{how} $!A$ was !!{prove}d. To this end, we can define a
\emph{term-labelled} system \Log{tN2i}, in which sequents have the form
$\Gamma \Sequent N : !A$. $\Gamma$~and $!A$ are as before: a set of
pairs of the form $x: !B$ and the !!{formula} the sub-!!{proof} ending
in the sequent has shown to follow from the assumptions in~$\Gamma$.
$N$~will be a certain term that encodes the !!{proof} ending in the
sequent.

The terms we assign to each sequent are built from variables $x$, $y$,
\dots, using function symbols which encode the inferences carried out
so far. We need a separate function symbol for each inference rule.
The function symbol corresponding to !!{introduction} rules are called
``constructors'' and those corresponding to !!{elimination} rules
``destructors.'' Each function symbol takes as many arguments as the
rule has premises. For instance, the function symbols for
\Intro{\land} and \Elim{\land}[1] might be $c^\land$ (2-place) and
$d^\land_1$ (1-place). In \Log{tN2i} these rules then become
\[
\Axiom$\Gamma_1 \fCenter N: !A$
\Axiom$\Gamma_2 \fCenter M: !B$
\RightLabel{\Intro{\land}}
\BinaryInf$\Gamma_1, \Gamma_2 \fCenter c^\land(N, m) : !A \land !B $
\DisplayProof
\quad
\Axiom$\Gamma \fCenter N: !A \land !B$
\RightLabel{\Elim{\land}[1]}
\UnaryInf$\Gamma \fCenter d^\land_1(N) : !A$
\DisplayProof
\]
If an inference rule discharges an assumption, i.e., the rule has a
premise that indicates a labelled !!{formula}~$x:!A$ in a premise that
is not in the context of the conclusion, we have to indicate this in
the proof term as well. The function symbols associate with such rules
\emph{bind} the corresponding variables. For instance, the
\Intro{\lif} rule goes from the premise $x:!A, \Gamma \Sequent !B$ to
$\Gamma \Sequent !A \lif !B$. The constructor $c^{\lif}$ takes one
argument, but also must indicate that the variable $x$ becomes bound.
This is the proof term's way of expressing that the assumption with
label~$x$ has been discharged. For instance, if the proof term for the
premise is~$N$, we could write the proof term for the conclusion as
$c^{\lif}([x]N)$ and the rule as
\[
\Axiom$x:!A, \Gamma \fCenter N: !B$
\RightLabel{\Intro{\lif}}
\UnaryInf$\Gamma \fCenter c^{\lif}([x]N) : !A \lif !B$
\DisplayProof
\]

In order to be able to fully reconstruct !!a{proof}, we need some
additional information. If the premise !!{formula}s of a rule don't
fully determine what the conclusion !!{formula} is, we have to add
that information to the term. This is the case in \Intro{\lif}, so we
will write $c^{\lif}([x:!A]N)$. It is also the case in \Intro{\lor}
and~\FalseInt, so we use function symbols that carry that information,
e.g.,
\[
\Axiom$\Gamma \fCenter N:!A$
\RightLabel{\Intro{\lor}[1]}
\UnaryInf$\Gamma \fCenter c^{\lor{!B}}_1:!A \lor !B$
\DisplayProof
\quad
\Axiom$\Gamma \fCenter N:\lfalse$
\RightLabel{\FalseInt}
\UnaryInf$\Gamma \fCenter d^\lfalse_{!A}(N):!A$
\DisplayProof
\]

With these ideas it is possible to construct a sequent-style natural
deduction system in which every sequent is of the form $\Gamma
\Sequent N : !A$, where $N$~is a \emph{proof term}.

It turns out that there is a close connection between such proof terms
and terms in the so-called \emph{typed lambda calculus}. To respect
this connection, we'll use symbols used in the literature on the
lambda calculus rather than the generic constructor and destructor
symbols above. E.g., instead of writing $c^{\lif}([x:!A]N)$ we write
$\lambd[\typeof{x}{!A}][N]$, instead of $d^{\lif}(N,M)$ we write $(NM)$,
$\pair{N}{M}$ instead of $c^\land(N,M)$, and $\proj{i}{N}$ instead
of~$d^\land_i(N)$.

\begin{defn}
  Proof terms are inductively generated by the following rules:
  \begin{enumerate}
  \item Every variable $x$ is a proof term (Axioms).
  \item If $x$ is a variable, $!A$ is !!a{formula}, and $N$ is a proof
    term, then $\lambd[\typeof{x}{!A}][N]$ is also a proof term~(\Intro\lif).
  \item If $N$ and $M$ are proof terms, then $(NM)$ is also a proof
    term~(\Elim\lif).
  \item If $N$ and $M$ are proof terms, then $\pair{N}{M}$ is a proof
    term~(\Intro\land).
  \item If $N$ is a proof term, then $\proj{i}{N}$ is also a proof
    term, where $i$ is $1$ or~$2$~(\Elim\land[i]).
  \item If $N$ is a proof term, and $!A$ is a formula, then
    $\inj[!A]{i}{!M}$ is a proof term, where $i$ is $1$ or~$2$~(\Intro\lor[i]).
  \item If $M$, $N_1$, $N_2$ are proof terms, and $x_1$, $x_2$ are
    variables, then $\dcase{M}{x_1}{N_1}{x_2}{N_2}$ is a proof
    term~(\Elim\lor).
  \item If $N$ is a proof term and $!A$ is !!a{formula}, then
    $\abort{!A}{N}$ is proof term~(\FalseInt).
  \end{enumerate}
\end{defn}

Not every proof term is the proof term of !!a{proof}. For instance,
a term $\pair{N}{M}$ encodes a proof ending in a~\Intro{\land}
inference, so its end-!!{formula} is a conjunction~$!A \land !B$. This
can't be the major premise of a~\Elim{\lif} rule, so $(\pair{N}{M}P)$
can't be the proof term of a correct proof. Only those proof terms
that follow the rules of \Log{tN2ip} are \emph{correct} proof terms.
These rules are given in \olref{tab:tN2ip}.

\subfile{rules-tN2}

\end{document}
